\documentclass[hyperref={pdfpagelabels=false},table]{beamer}

\input{lec_common}
\begin{document}

\part{Unit 6 (Cont.): Object oriented programming with classes}

%\frame{\tableofcontents}

%%%%%%%%%%%%%%%%%%%%%%%%%%%%%%%%%%%%%%%%%%%%%%%%%%%%%%%%%%%%%%%%%%%%%
%%%%%%%%%%%%%%%%%%%%%%%%%%%%%%%%%%%%%%%%%%%%%%%%%%%%%%%%%%%%%%%%%%%%%
\begin{frame}[fragile]\frametitle{Example: RationalNumber}

  Let's reconsider the example with the class \pyth{RationalNumber}.

\begin{python}
class RationalNumber:
    def __init__(self,numerator,denominator):
        if not isinstance(numerator,int) :
             raise TypeError(
                   "numerator should be of type int")
        if not isinstance(denominator,int) :
             raise TypeError(
                   "denominator should be of type int")
        self.numerator=numerator
        self.denominator=denominator
\end{python}

\end{frame}


%%%%%%%%%%%%%%%%%%%%%%%%%%%%%%%%%%%%%%%%%%%%%%%%%%%%%%%%%%%%%%%%%%%%%
%%%%%%%%%%%%%%%%%%%%%%%%%%%%%%%%%%%%%%%%%%%%%%%%%%%%%%%%%%%%%%%%%%%%%
\begin{frame}[fragile]\frametitle{Everything can be changed}

  By \alert{default} all members in a Python class are \alert{public}. 
  Any member can be accessed from outside the class environment. 
  Everything can be changed after creating an instance. \\ \ \\

  In other programming languages such as C++ or Java there exist 
  concepts for public, protected and private members in a class. \\ \ \\

  \pause

\begin{python}
  r = RationalNumber(2,3)
  print(r)                        # 2/3 
  r.numerator = 23                # error-prone  
  r.name = "my favorite fraction" # a new attribute! 
  print(r)                        # 23/3
\end{python}

\pause

We can even delete an attribute!

\begin{python}
  r = RationalNumber(2,3)
  del r.numerator   
  print(r)                        # error
\end{python}

\end{frame}

%%%%%%%%%%%%%%%%%%%%%%%%%%%%%%%%%%%%%%%%%%%%%%%%%%%%%%%%%%%%%%%%%%%%%
%%%%%%%%%%%%%%%%%%%%%%%%%%%%%%%%%%%%%%%%%%%%%%%%%%%%%%%%%%%%%%%%%%%%%
\begin{frame}[fragile]\frametitle{Changing attributes}

We can \alert{protect} attributes. 
\structure{By convention} attributes that have a leading underscore 
are meant for internal use only.

\begin{python}
class RationalNumber:
    def __init__(self,numerator,denominator):
        ....
        self._numerator=numerator
        self._denominator=denominator
    def getNum(self):
        return self._numerator
    def setNum(self, numerator):
        self._numerator = numerator
\end{python}

\pause 

\begin{python}
r = RationalNumber(2, 3)
print(r._numerator)  # bad style
print(r.getNum())    # good style
r._numerator = 23    # bad style
r.setNum( 23 )       # good style
\end{python}

\end{frame}

%%%%%%%%%%%%%%%%%%%%%%%%%%%%%%%%%%%%%%%%%%%%%%%%%%%%%%%%%%%%%%%%%%%%%
%%%%%%%%%%%%%%%%%%%%%%%%%%%%%%%%%%%%%%%%%%%%%%%%%%%%%%%%%%%%%%%%%%%%%
\begin{frame}\frametitle{A triangle example}

  \includegraphics[width=\textwidth]{triangle}

\end{frame}


\begin{frame}[fragile]{Attributes depending on each other -- A guiding example}
Let us assume the following class to define a triangle
\begin{python}
class Triangle:
    def __init__(self, A, B, C):
        self.A = array(A)
        self.B = array(B)
        self.C = array(C)
        self.a = self.C - self.B
        self.b = self.C - self.A
        self.c = self.B - self.A
    def area(self):
        return abs(cross(self.b, self.c))/2
\end{python}
An instance of this triangle is created by
\begin{python}
tr = Triangle([0.,0.], [1.,0.], [0.,1.])
\end{python}
and its area is computed by
\begin{python}
tr.area() # returns 0.5
\end{python}
\end{frame}
\begin{frame}[fragile]{Altering an attribute}
Altering an attribute,
\begin{python}
tr.B = array([1.0, 0.1])
\end{python}
does not affect the attribute 'area':
\begin{python}
tr.area() # still 0.5
\end{python}
We have to be able to run ``internally'' some commands when one attribute is
changed in order to modify the others!
\vfill
For this end we introduce attributes only for internal use, e.g \pyth!_B! \\
and ...
\end{frame}

%%%%%%%%%%%%%%%%%%%%%%%%%%%%%%%%%%%%%%%%%%%%%%%%%%%%%%%%%%%%%%%%%%%%%
%%%%%%%%%%%%%%%%%%%%%%%%%%%%%%%%%%%%%%%%%%%%%%%%%%%%%%%%%%%%%%%%%%%%%
\begin{frame}\frametitle{The property command}

We can bind the getter and setter methods to a property \pyth{B}. 
The property B is accessed directly but internally the getter and setter 
methods will use the hidden attribute \pyth{_B}.

\end{frame}

%%%%%%%%%%%%%%%%%%%%%%%%%%%%%%%%%%%%%%%%%%%%%%%%%%%%%%%%%%%%%%%%%%%%%
%%%%%%%%%%%%%%%%%%%%%%%%%%%%%%%%%%%%%%%%%%%%%%%%%%%%%%%%%%%%%%%%%%%%%

\begin{frame}[fragile]{Altering an attribute \hfill (Cont.)}
... and modify our class
\begin{python}
class Triangle:
    def __init__(self,A, B, C):
        self._A = array(A)
        self._B = array(B)
        self._C = array(C)
        self._a = self._C - self._B
        self._b = self._C - self._A
        self._c = self._B - self._A
    def area(self):
        return abs(cross(self._c, self._b)) / 2.
    def setB(self,B):
        self._B = B
        self._a = self._C - self._B
        self._c = self._B - self._A
    def getB(self):
        return self._B
    B=property(fget = getB, fset = setB)
\end{python}
\vfill
{\small What would happen if we change \pyth!self._B! to \pyth!self.B!?}
\end{frame}
\begin{frame}[fragile]{A triangle has three corners}
Python allows us to delete an attribute
\vfill
\begin{python}
del tr.B
\end{python}
If so, we won't have a triangle any more. So let's prohibit deletion
\begin{python}
class Triangle:
  ...
  def delPoint(self):
    raise Exception("A triangle point cannot be deleted")
  B = property(fget = getB,
               fset = setB,
               fdel = delPoint)
\end{python}
\vfill
\begin{python}
del tr.B # now raises an exception
\end{python}

\end{frame}

\begin{frame}[fragile]{Inheritance - a mathematical concept}
Mathematics builds on class inheritance principles:

\vfill

Example of a mathematical inheritance hierarchy

\begin{itemize}
\item {\em Functions} \\ {\small can be evaluated, differentiated, integrated, summed up, ....}
% \item {\em Polynomials} are functions\\ {\small we can delete coefficients, ...}
% \item {\em Trigonometric polynomials} are polynomials \\ {\small same methods as polynomials, change representation, ask for real parts, truncate, ...}
\item {\em Polynomials} are functions\\ {\small same methods as functions, change representation, truncate, ...}
\item {\em Trigonometric polynomials} are polynomials \\ {\small same methods as polynomials, ask for real parts, ...}
\end{itemize}
\vfill

We say {\it polynomials inherit properties and methods from functions}.
\end{frame}

%%%%%%%%%%%%%%%%%%%%%%%%%%%%%%%%%%%%%%%%%%%%%%%%%%%%%%%%%%%%%%%%%%%%%
%%%%%%%%%%%%%%%%%%%%%%%%%%%%%%%%%%%%%%%%%%%%%%%%%%%%%%%%%%%%%%%%%%%%%
\begin{frame}[fragile]\frametitle{Inheritance}

Example -- lines

A line can be defined by two points $A$ and $B$.

We will define a parent class, or base class, \pyth{Line}.

  Then we will define two child classes, or derived classes, \pyth{DashedLine} and \pyth{NonVerticalLine}.

To make a class that is a child class of som parent class, the parent class name is written inside parenthesis:

\begin{python}
class ChildClass(ParentClass):
    ...
\end{python}    

\end{frame}
%%%%%%%%%%%%%%%%%%%%%%%%%%%%%%%%%%%%%%%%%%%%%%%%%%%%%%%%%%%%%%%%%%%%%
%%%%%%%%%%%%%%%%%%%%%%%%%%%%%%%%%%%%%%%%%%%%%%%%%%%%%%%%%%%%%%%%%%%%%
\begin{frame}[fragile]\frametitle{The base class}

\begin{python}
class Line:
    def __init__(self, A, B):
        self.A = array(A)
        self.B = array(B) 
    def plot(self):
        grid()  #  grid and plot from pyplot
        plot([self.A[0], self.B[0]], [self.A[1], self.B[1]])

L = Line([1, 1], [5, 3])
L.plot()
\end{python}

\end{frame}

%%%%%%%%%%%%%%%%%%%%%%%%%%%%%%%%%%%%%%%%%%%%%%%%%%%%%%%%%%%%%%%%%%%%%
%%%%%%%%%%%%%%%%%%%%%%%%%%%%%%%%%%%%%%%%%%%%%%%%%%%%%%%%%%%%%%%%%%%%%
\begin{frame}[fragile]\frametitle{Defining a derived class \pyth{DashedLine}}

  A \pyth{DashedLine} inherits all data attributes and all methods 
  from \pyth{Line} but the plot method is redefined, 
  overriding the method of the base class.

\begin{python}
class DashedLine(Line): # indicate that class 
                        # is derived from Line
    def plot(self):  # override the plot method
        plot([self.A[0], self.B[0]], [self.A[1], self.B[1]],"--")

dl = DashedLine([1, 1], [5, 3])
print(f"dl.A = {dl.A}, dl.B = {dl.B}")
dl.plot() 
\end{python}

\end{frame}

%%%%%%%%%%%%%%%%%%%%%%%%%%%%%%%%%%%%%%%%%%%%%%%%%%%%%%%%%%%%%%%%%%%%%
%%%%%%%%%%%%%%%%%%%%%%%%%%%%%%%%%%%%%%%%%%%%%%%%%%%%%%%%%%%%%%%%%%%%%
\begin{frame}[fragile]\frametitle{Using \pyth{super()}}

  We can make another derived \pyth{class NonVerticalLine}. 
  It has the same attributes as \pyth{Line}, 
  but also the attributes \pyth{slope} and \pyth{y_intercept}.

  To call the \pyth{__init__} method of the parent class, 
  we refer to the parent class by using \pyth{super()}.

\begin{python}
class NonVerticalLine(Line):
    def __init__(self, A, B): # override init method
        super().__init__(A, B)
        self.slope = (B[1]-A[1])/(B[0]-A[0]) # may raise a ZeroDivisionError
        self.yIntercept = A[1] - self.slope*A[0]

    def f(self, x):
        return self.slope*x + self.yIntercept

L = NonVerticalLine([-1, -1], [4, 1])
print(f"L.A = {L.A}, L.B = {L.B}, L.slope = {L.slope}, L.yIntercept = {L.yIntercept}")
L.plot()
\end{python}

\end{frame}

%%%%%%%%%%%%%%%%%%%%%%%%%%%%%%%%%%%%%%%%%%%%%%%%%%%%%%%%%%%%%%%%%%%%%
%%%%%%%%%%%%%%%%%%%%%%%%%%%%%%%%%%%%%%%%%%%%%%%%%%%%%%%%%%%%%%%%%%%%%
\begin{frame}[fragile]\frametitle{\pyth{isinstance} and \pyth{issubclass}}

A non-vertical line is also a line.

\begin{python}
L = NonVerticalLine([1, 2], [2, 4])

print(isinstance(L, Line))
print(isinstance(L, NonVerticalLine))
print(isinstance(L, DashedLine))
\end{python}

\pause 

We can check if a class is a derived class of another class

\begin{python}
# NonVerticalLine is a subclass of Line
print(issubclass(NonVerticalLine, Line))  
# Line is not a subclass of NonVerticalLine
print(issubclass(Line, NonVerticalLine)) 
\end{python}

\end{frame}

%%%%%%%%%%%%%%%%%%%%%%%%%%%%%%%%%%%%%%%%%%%%%%%%%%%%%%%%%%%%%%%%%%%%%
%%%%%%%%%%%%%%%%%%%%%%%%%%%%%%%%%%%%%%%%%%%%%%%%%%%%%%%%%%%%%%%%%%%%%
\begin{frame}[fragile]\frametitle{Objects that act like functions}

Recall that the \pyth{quad} function can only be used on functions with one single parameter. 

\begin{python}
from scipy.integrate import quad

def f(x):
    return 2*sin(0.5*x)

def g(x, A, omega):
    return A*sin(omega*x)

a = quad(f, 0, pi) # works
print(a)

b = quad(g, 0, pi) # doesn't work
\end{python}

\end{frame}

%%%%%%%%%%%%%%%%%%%%%%%%%%%%%%%%%%%%%%%%%%%%%%%%%%%%%%%%%%%%%%%%%%%%%
%%%%%%%%%%%%%%%%%%%%%%%%%%%%%%%%%%%%%%%%%%%%%%%%%%%%%%%%%%%%%%%%%%%%%
\begin{frame}[fragile]\frametitle{Objects that act like functions (cont.)}

We can make a class that has a value method corresponding to a function value.  

\begin{python}
class Sin:
    def __init__(self, A, omega):
        self.A = A
        self.omega = omega
    def value(self, x):
        return self.A*sin(self.omega*x)
    def __repr__(self):
        return f"{self.A}*sin({self.omega}x)"

sin1 = Sin(1, 1)
sin2 = Sin(2, 0.5)
print(sin1)
print(sin2)

a = quad(sin1.value, 0, pi)
b = quad(sin2.value, 0, pi)
print(a)
print(b)
\end{python}

\end{frame}

%%%%%%%%%%%%%%%%%%%%%%%%%%%%%%%%%%%%%%%%%%%%%%%%%%%%%%%%%%%%%%%%%%%%%
%%%%%%%%%%%%%%%%%%%%%%%%%%%%%%%%%%%%%%%%%%%%%%%%%%%%%%%%%%%%%%%%%%%%%
\begin{frame}[fragile]\frametitle{The special method \pyth{__call__}}

\begin{python}
class Sin:
    def __init__(self, A, omega):
        self.A = A
        self.omega = omega
    def __call__(self, x): # easier to use than the
                           # previous value method 
        return self.A*sin(self.omega*x)
    def __repr__(self):
        return f"{self.A}*sin({self.omega}x)"

sin1 = Sin(1, 1)
sin2 = Sin(2, 0.5)
print(sin1(pi/2))  # print sin at position pi/2, 
print(sin2(pi))    # here __call__ is executed

a = quad(sin1, 0, pi) # sin1.__call__ is used
b = quad(sin2, 0, pi) # sin2.__call__ is used
print(a)
print(b)
\end{python}  
\end{frame}

%%%%%%%%%%%%%%%%%%%%%%%%%%%%%%%%%%%%%%%%%%%%%%%%%%%%%%%%%%%%%%%%%%%%%
%%%%%%%%%%%%%%%%%%%%%%%%%%%%%%%%%%%%%%%%%%%%%%%%%%%%%%%%%%%%%%%%%%%%%
\begin{frame}[fragile]\frametitle{Arithmetic operators for functions}
In mathematics we can add functions.

 \begin{eqnarray*} 
   f(x) &=& sin(x) \\
   g(x) &=& cos(x) \\
   h(x) &=& f(x) + g(x)
 \end{eqnarray*}

\pause 

In Python this is not directly possible:

\begin{python}
f = sin
g = cos
h = f + g  # error 
\end{python}

\end{frame}

%%%%%%%%%%%%%%%%%%%%%%%%%%%%%%%%%%%%%%%%%%%%%%%%%%%%%%%%%%%%%%%%%%%%%
%%%%%%%%%%%%%%%%%%%%%%%%%%%%%%%%%%%%%%%%%%%%%%%%%%%%%%%%%%%%%%%%%%%%%
\begin{frame}[fragile]\frametitle{Adding functions}

With the concept of classes, we can program an elegant work-around.

\begin{python}
class Function: # encapsulate a function
    def __init__(self, f):
        self.f = f
    def __call__(self, x):
        return self.f(x)
    def __add__(self, g):
        def sum(x):
            return self(x)+g(x)
        return type(self)(sum) # return an instance of the same class

f = Function(sin)
g = Function(cos)
h = f + g
x = linspace(-pi, pi)
plot(x, f(x), "--", x, g(x), "--", x, h(x))
\end{python}

\end{frame}


%%%%%%%%%%%%%%%%%%%%%%%%%%%%%%%%%%%%%%%%%%%%%%%%%%%%%%%%%%%%%%%%%%%%%
%%%%%%%%%%%%%%%%%%%%%%%%%%%%%%%%%%%%%%%%%%%%%%%%%%%%%%%%%%%%%%%%%%%%%
\begin{frame}[fragile]\frametitle{Class attributes, static and class methods}

  In following class, \pyth{_clsVar} is a class data attribute, 
  \pyth{_instVar} is an instance data attribute.

\begin{python}
class A:
  _clsVar = 0 # class attribute, or static attribute
  
  def __init__(self):
    self._instVar = 0 # an instance data attribute

  def instMth(self, x): # the first arg is the inst
    print(f"executing instMth({self}, {x})")
        
  @classmethod
  def classMth(cls, x): # the first arg is the class
    print(f"executing classMth({cls}, {x})")

  @staticmethod
  def staticMth(x): # a function bound to the class
    print(f"executing staticMth({x})")
\end{python}  

\end{frame}

%%%%%%%%%%%%%%%%%%%%%%%%%%%%%%%%%%%%%%%%%%%%%%%%%%%%%%%%%%%%%%%%%%%%%
%%%%%%%%%%%%%%%%%%%%%%%%%%%%%%%%%%%%%%%%%%%%%%%%%%%%%%%%%%%%%%%%%%%%%
\begin{frame}[fragile]\frametitle{Class attributes, static and class methods}

\begin{python}
a = A()
# What happens?
a.instMth(1)
A.instMth(1)
a.classMth(2)
A.classMth(2)
a.staticMth(3)
A.staticMth(3)
\end{python}

\end{frame}

%%%%%%%%%%%%%%%%%%%%%%%%%%%%%%%%%%%%%%%%%%%%%%%%%%%%%%%%%%%%%%%%%%%%%
%%%%%%%%%%%%%%%%%%%%%%%%%%%%%%%%%%%%%%%%%%%%%%%%%%%%%%%%%%%%%%%%%%%%%
\begin{frame}[fragile]\frametitle{Class attributes -- \pyth{refcounter.py}}

\begin{python}
class RationalNumber:
  _nInstances = 0 # class attribute
  def __init__(self ,numerator ,denominator):
    self._numerator   = numerator
    self._denominator = denominator

    # increase reference counter, note the self here
    self.inc()

  # destructor called when object is deleted
  def __del__(self):
    # decrease reference counter
    self.dec()

  @classmethod # increase class attribute by one
  def inc(cls):
    cls._nInstances += 1

  @classmethod
  def dec(cls): # decrease class attribute by one
    cls._nInstances -= 1
\end{python}

\end{frame}

%%%%%%%%%%%%%%%%%%%%%%%%%%%%%%%%%%%%%%%%%%%%%%%%%%%%%%%%%%%%%%%%%%%%%
%%%%%%%%%%%%%%%%%%%%%%%%%%%%%%%%%%%%%%%%%%%%%%%%%%%%%%%%%%%%%%%%%%%%%
\begin{frame}\frametitle{Summary}

  \begin{itemize}
    \item Instance methods needs an object and can access the object through a parameter \pyth{self}.
    \item Class methods are in the namespace of a class and can access the class through a parameter \pyth{cls}.
    \item Static methods are in the namespace of a class but cannot access either the class or an instance.
  \end{itemize}

  \pause

  \vskip 5pt 

  The words \pyth{self} and \pyth{cls} are not reserved words but conventions.
\end{frame}


%%%%%%%%%%%%%%%%%%%%%%%%%%%%%%%%%%%%%%%%%%%%%%%%%%%%%%%%%%%%%%%%%%%%%
%%%%%%%%%%%%%%%%%%%%%%%%%%%%%%%%%%%%%%%%%%%%%%%%%%%%%%%%%%%%%%%%%%%%%

\begin{comment}
\begin{frame}[fragile]{Inheritance in Python: An example}
\vfill
To solve an ordinary differential equation
\[
\dot x = f(t,x) \;\;x(0)=x_0
\]
you used in another course Euler's method
\[
 u_{i+1} = u_i + h f(t_i, u_i), \;\;\; \text{with } u_0=x_0\;\text{and}\;t_i=t_{i-1}+h .
\]
\vfill
It is a special case of a family of methods of the type
\[
 u_{i+1} = u_i + h \Phi(f,u_i,t_i,h) .
\]
\end{frame}
\begin{frame}[fragile]{Inheritance - in Python}
We consider a general integrator class performing:
\[
 u_{i+1} = u_i + h \Phi(f,u_i,t_i,h) .
\]
\begin{python}
class Onestepmethod:
  def __init__(self, f, x0, t0, te, N):
    self.f = f
    self.x0 = x0
    self.interval = [t0, te]
    self.grid = linspace(t0, te, N)
    self.h = (te - t0)/N
  def step(self):
    ti, ui = self.grid[0], self.x0
    yield ti, ui
    for t in self.grid[1:]:
      ui = ui + self.h*self.phi(self.f,ui,ti)
      ti = t
      yield ti, ui        # see one of the next lectures
  def solve(self):
    self.solution = array([tu for tu in self.step()])
\end{python}

\end{frame}
\begin{frame}[fragile]{Inheritance - in Python \hfill (cont.)}
 ... construct from it two subclasses
 \begin{python}
class ExpEuler(Onestepmethod):
    def phi(self, f, u, t):
        return f(u, t)

class MidPointRule(Onestepmethod):
    def phi(self, f, u, t):
        return f(u + self.h/2*f(u,t), t + self.h/2)
\end{python}

\end{frame}
\begin{frame}[fragile]{Inheritance - in Python \hfill (cont.)}
Which we then call
\begin{python}
def f(x, t):
    return -0.5*x

euler = ExpEuler(f, 15., 0., 10., 20)
euler.solve()

midpoint = MidPointRule(f, 15., 0., 10., 20)

midpoint.solve()
\end{python}
\end{frame}
\end{comment}

\end{document}


